\documentclass[11pt]{article}

\usepackage[margin=1in]{geometry}
\usepackage{amsmath, amssymb, amsthm}
\usepackage{mathtools}
\usepackage{bm}
\usepackage{hyperref}
\usepackage{cleveref}
\usepackage{enumitem}
\usepackage{booktabs}
\usepackage{array}
\usepackage[T1]{fontenc}
\usepackage{lmodern}

% ── Theorem environments ──────────────────────────────────────────────────────
\newtheorem{theorem}{Theorem}[section]
\newtheorem{proposition}[theorem]{Proposition}
\newtheorem{definition}[theorem]{Definition}
\newtheorem{remark}[theorem]{Remark}

% ── Notation shortcuts ────────────────────────────────────────────────────────
\newcommand{\R}{\mathbb{R}}
\newcommand{\dd}{\,\mathrm{d}}
\newcommand{\pderiv}[2]{\frac{\partial #1}{\partial #2}}
\newcommand{\ddt}{\frac{\mathrm{d}}{\mathrm{d}t}}
\newcommand{\bq}{\mathbf{q}}
\newcommand{\bp}{\mathbf{p}}
\newcommand{\bv}{\mathbf{v}}
\newcommand{\bF}{\mathbf{F}}
\newcommand{\J}{\mathbf{J}}          % Jacobian
\newcommand{\Omg}{\boldsymbol{\Omega}} % symplectic matrix

\hypersetup{
  colorlinks=true,
  linkcolor=blue,
  citecolor=blue,
  urlcolor=blue,
  pdftitle={Symplectic Euler Integration},
  pdfauthor={Ernest Yeung},
}

% ─────────────────────────────────────────────────────────────────────────────
\title{%
  \textbf{Symplectic Euler Integration}\\[4pt]
  \large From Hamiltonian Mechanics to Real-Time Physics Simulation
}
\author{Ernest Yeung}
\date{June 2026}

\begin{document}
\maketitle
\tableofcontents
\bigskip

% ─────────────────────────────────────────────────────────────────────────────
\section{Introduction}

Physics engines such as NVIDIA PhysX (used inside Isaac Sim) must integrate
Newton's equations of motion millions of times per second at a fixed timestep
$h$.  Naive application of the simplest numerical integrator---explicit
(forward) Euler---causes energy to grow without bound, producing bodies that
accelerate spontaneously.  The standard remedy is the \emph{symplectic Euler}
(also called \emph{semi-implicit Euler}) method, which is the unique first-order
scheme that exactly preserves the geometric structure underlying Hamiltonian
mechanics.

This note derives that structure from first principles, proves that explicit
Euler breaks it while symplectic Euler preserves it, connects the method to
the classical generating-function theory of canonical transformations, and works
out the exact discrete free-fall formula that emerges---the formula used to
validate the PhysX integrator in the Starship IsaacSim test suite (see
\texttt{tests/test\_t1\_free\_fall.py}).

% ─────────────────────────────────────────────────────────────────────────────
\section{Hamiltonian Mechanics and the Symplectic Form}

\subsection{Phase Space and Hamilton's Equations}

Let $q \in \R^n$ denote generalised coordinates and $p \in \R^n$ the
conjugate momenta.  The pair $(q, p)$ lives in the \emph{phase space}
$T^*Q \cong \R^{2n}$.  Given a Hamiltonian $H : T^*Q \to \R$, the equations
of motion are
\begin{equation}\label{eq:hamilton}
  \dot{q} = \pderiv{H}{p}, \qquad
  \dot{p} = -\pderiv{H}{q}.
\end{equation}
These can be written compactly as $\dot{z} = \Omg^{-1} \nabla H(z)$, where
$z = (q, p)^{\top}$ and
\begin{equation}\label{eq:symplectic_matrix}
  \Omg =
  \begin{pmatrix} 0 & I_n \\ -I_n & 0 \end{pmatrix}
  \in \R^{2n \times 2n}
\end{equation}
is the \emph{standard symplectic matrix} (satisfying
$\Omg^{\top} = -\Omg$, $\Omg^2 = -I$).

\subsection{The Symplectic 2-Form}

On $T^*Q$ there is a canonical closed, non-degenerate 2-form
\begin{equation}\label{eq:2form}
  \omega = \sum_{i=1}^n \dd q_i \wedge \dd p_i.
\end{equation}
In coordinates, a linear map $\Phi : \R^{2n} \to \R^{2n}$ with matrix $M$
is called \emph{symplectic} if it preserves $\omega$:
\begin{equation}\label{eq:symplectic_condition}
  M^{\top} \Omg\, M = \Omg.
\end{equation}
For $n=1$ this reduces to $\det(M) = 1$, i.e., the map preserves oriented
area in the $(q, p)$-plane.

\subsection{Liouville's Theorem}

\begin{theorem}[Liouville]
  The flow $\varphi_t$ of a Hamiltonian system \eqref{eq:hamilton} is
  symplectic for all $t$: $(\varphi_t)^* \omega = \omega$.  In particular it
  preserves the $2n$-dimensional phase-space volume $\dd q\, \dd p$.
\end{theorem}

\begin{proof}[Sketch]
  Differentiating $M(t) = D\varphi_t(z_0)$ along the flow and using
  $\dot{M} = JM$ where $J = \Omg^{-1} D^2 H$ is the linearised vector field
  gives $\frac{\dd}{\dd t}(M^{\top}\Omg M) = 0$, so $M^{\top}\Omg M$ is
  constant.  At $t=0$, $M(0) = I$ and $I^{\top}\Omg I = \Omg$, proving the
  claim for all $t$.
\end{proof}

% ─────────────────────────────────────────────────────────────────────────────
\section{Explicit Euler: Not Symplectic}

\subsection{The Algorithm}

For a particle of mass $m$ with Hamiltonian
$H = \tfrac{p^2}{2m} + V(q)$ (one degree of freedom, $n=1$), Hamilton's
equations are
\begin{equation}
  \dot{q} = \frac{p}{m}, \qquad \dot{p} = F(q) \coloneqq -V'(q).
\end{equation}
The \textbf{explicit (forward) Euler} discretisation uses state values from
step $n$ to advance to step $n+1$:
\begin{align}
  p_{n+1} &= p_n + h\, F(q_n), \label{eq:euler_p}\\
  q_{n+1} &= q_n + \frac{h}{m}\, p_n. \label{eq:euler_q}
\end{align}
Both updates use the \emph{old} values $(q_n, p_n)$.

\subsection{Jacobian and Failure of Symplecticity}

The Jacobian of the map $(q_n, p_n) \mapsto (q_{n+1}, p_{n+1})$ is
\begin{equation}\label{eq:euler_jacobian}
  \J_{\mathrm{expl}} =
  \begin{pmatrix}
    \partial q_{n+1}/\partial q_n & \partial q_{n+1}/\partial p_n \\
    \partial p_{n+1}/\partial q_n & \partial p_{n+1}/\partial p_n
  \end{pmatrix}
  =
  \begin{pmatrix}
    1 & h/m \\
    h F'(q_n) & 1
  \end{pmatrix}.
\end{equation}
Its determinant is
\begin{equation}\label{eq:euler_det}
  \det(\J_{\mathrm{expl}}) = 1 - \frac{h^2}{m} F'(q_n).
\end{equation}
For any non-trivial potential ($F' \neq 0$) this differs from 1, so explicit
Euler \emph{does not} satisfy the symplectic condition
\eqref{eq:symplectic_condition}.

\paragraph{Consequence.}
For a harmonic oscillator with $F(q) = -kq$ (spring constant $k > 0$),
$F'(q) = -k$, giving
\[
  \det(\J_{\mathrm{expl}}) = 1 + \frac{h^2 k}{m} > 1.
\]
Each step \emph{expands} phase-space area by the factor
$1 + h^2 k/m > 1$, which compounds over $N$ steps to
$(1 + h^2 k/m)^N \to \infty$.  This corresponds to spurious energy growth
that makes explicit Euler numerically unstable for oscillatory systems at
any fixed $h > 0$.

% ─────────────────────────────────────────────────────────────────────────────
\section{Symplectic Euler: Preserving the Structure}

\subsection{The Algorithm}

The \textbf{symplectic (semi-implicit) Euler} method updates momentum first
using the old position, then updates position using the \emph{new} momentum:
\begin{align}
  p_{n+1} &= p_n + h\, F(q_n), \label{eq:symp_p}\\
  q_{n+1} &= q_n + \frac{h}{m}\, p_{n+1}. \label{eq:symp_q}
\end{align}
The only difference from \eqref{eq:euler_p}--\eqref{eq:euler_q} is that
\eqref{eq:symp_q} uses $p_{n+1}$ rather than $p_n$.

\subsection{Jacobian and Proof of Symplecticity}

Substituting \eqref{eq:symp_p} into \eqref{eq:symp_q}:
\[
  q_{n+1} = q_n + \frac{h}{m}(p_n + hF(q_n))
           = q_n + \frac{h}{m}p_n + \frac{h^2}{m}F(q_n).
\]
The Jacobian is therefore
\begin{equation}\label{eq:symp_jacobian}
  \J_{\mathrm{symp}} =
  \begin{pmatrix}
    1 + \dfrac{h^2}{m}F'(q_n) & \dfrac{h}{m} \\[8pt]
    h F'(q_n) & 1
  \end{pmatrix}.
\end{equation}

\begin{proposition}\label{prop:symplectic}
  The map \eqref{eq:symp_p}--\eqref{eq:symp_q} is symplectic, i.e.,
  $\J_{\mathrm{symp}}^{\top}\Omg\,\J_{\mathrm{symp}} = \Omg$ (equivalently,
  $\det(\J_{\mathrm{symp}}) = 1$).
\end{proposition}

\begin{proof}
  Let $\alpha = h^2 F'(q_n)/m$, $\beta = h/m$, $\gamma = h F'(q_n)$,
  so $\alpha = \beta\gamma$.  Then
  \[
    \J_{\mathrm{symp}} =
    \begin{pmatrix} 1+\alpha & \beta \\ \gamma & 1 \end{pmatrix}.
  \]
  We verify the symplectic identity with $\Omg = \bigl(\begin{smallmatrix}0 &
  1 \\ -1 & 0\end{smallmatrix}\bigr)$:
  \[
    \J^{\top}\Omg\J
    = \begin{pmatrix} 1+\alpha & \gamma \\ \beta & 1 \end{pmatrix}
      \begin{pmatrix} 0 & 1 \\ -1 & 0 \end{pmatrix}
      \begin{pmatrix} 1+\alpha & \beta \\ \gamma & 1 \end{pmatrix}.
  \]
  First factor on the right:
  \[
    \begin{pmatrix} 1+\alpha & \gamma \\ \beta & 1 \end{pmatrix}
    \begin{pmatrix} 0 & 1 \\ -1 & 0 \end{pmatrix}
    =
    \begin{pmatrix} -\gamma & 1+\alpha \\ -1 & \beta \end{pmatrix}.
  \]
  Then:
  \begin{align*}
    \J^{\top}\Omg\J
    &=
    \begin{pmatrix} -\gamma & 1+\alpha \\ -1 & \beta \end{pmatrix}
    \begin{pmatrix} 1+\alpha & \beta \\ \gamma & 1 \end{pmatrix} \\
    &=
    \begin{pmatrix}
      -\gamma(1+\alpha) + (1+\alpha)\gamma & -\gamma\beta + (1+\alpha) \\
      -(1+\alpha) + \beta\gamma & -\beta + \beta
    \end{pmatrix} \\
    &=
    \begin{pmatrix}
      0 & 1 + \alpha - \gamma\beta \\
      \beta\gamma - 1 - \alpha & 0
    \end{pmatrix}.
  \end{align*}
  Since $\alpha = \beta\gamma$:
  \[
    1 + \alpha - \gamma\beta = 1 + \beta\gamma - \gamma\beta = 1,\quad
    \beta\gamma - 1 - \alpha = \beta\gamma - 1 - \beta\gamma = -1.
  \]
  Therefore $\J^{\top}\Omg\J = \bigl(\begin{smallmatrix}0 & 1 \\ -1 &
  0\end{smallmatrix}\bigr) = \Omg$. \qed
\end{proof}

\noindent The key cancellation $\alpha = \beta\gamma$ is the algebraic
fingerprint of using $p_{n+1}$ in the position update: it makes the
off-diagonal contributions to $\det$ cancel exactly.

% ─────────────────────────────────────────────────────────────────────────────
\section{Derivation from the Generating-Function Theory}

\subsection{Type-2 Generating Functions}

Canonical transformation theory \cite{Goldstein2002} classifies
symplectomorphisms by the type of mixed-variable generating function
$S_k$ used to define them.  A \emph{Type-2 generating function}
$S_2(q, P)$ of old position $q$ and new momentum $P$ generates the
transformation via
\begin{equation}\label{eq:F2_relations}
  p = \pderiv{S_2}{q},\qquad Q = \pderiv{S_2}{P},
\end{equation}
and any $S_2$ satisfying these relations produces a canonical (symplectic)
map $(q, p) \mapsto (Q, P)$ by construction \cite{Hairer2006}.

\subsection{Symplectic Euler from a Type-2 Function}

Choose the generating function
\begin{equation}\label{eq:S2}
  S_2(q_n, p_{n+1}) = q_n \cdot p_{n+1} + h\, H(q_n, p_{n+1}).
\end{equation}
For $H = p^2/(2m) + V(q)$ this becomes
\[
  S_2(q_n, p_{n+1}) = q_n\, p_{n+1} + h\!\left[\frac{p_{n+1}^2}{2m} + V(q_n)\right].
\]
Applying \eqref{eq:F2_relations} with $q \leftarrow q_n$, $P \leftarrow p_{n+1}$,
$p \leftarrow p_n$, $Q \leftarrow q_{n+1}$:
\begin{align}
  p_n &= \pderiv{S_2}{q_n} = p_{n+1} + h\, V'(q_n)
       = p_{n+1} - h\, F(q_n),
  \label{eq:gen_p}\\
  q_{n+1} &= \pderiv{S_2}{p_{n+1}} = q_n + \frac{h}{m}\, p_{n+1}.
  \label{eq:gen_q}
\end{align}
Solving \eqref{eq:gen_p} for $p_{n+1}$:
\[
  p_{n+1} = p_n + h\, F(q_n),
\]
which is exactly \eqref{eq:symp_p}.  Substituting into \eqref{eq:gen_q}
recovers \eqref{eq:symp_q}.  Since every transformation generated by a
Type-2 function is canonical (it satisfies \eqref{eq:symplectic_condition}
by construction of the generating-function formalism), this provides an
independent proof of \cref{prop:symplectic} and reveals the \emph{origin}
of the method: symplectic Euler is the degree-1 truncation in $h$ of the
exact flow map, approximated via $S_2$.

\begin{remark}
  The identity $S_2(q, P) = qP$ generates the identity transformation.
  Adding $hH(q,P)$ as a correction gives a first-order approximation to the
  exact flow, but one that \emph{exactly} preserves the symplectic structure
  at every $h > 0$.
\end{remark}

% ─────────────────────────────────────────────────────────────────────────────
\section{Modified Hamiltonian and Bounded Energy Error}

A symplectic integrator does not conserve $H$ exactly; instead it exactly
conserves a \emph{modified Hamiltonian} $\tilde H$ close to $H$
\cite{Hairer2006,Leimkuhler2004}:
\begin{equation}\label{eq:modified_H}
  \tilde H(q,p) = H(q,p) + \frac{h}{2m}\,p\,V'(q) + \mathcal{O}(h^2).
\end{equation}
This is the key distinction from non-symplectic methods:
\begin{itemize}[leftmargin=2em]
  \item \textbf{Explicit Euler}: energy error grows \emph{monotonically} as
    $(1+h^2 k/m)^N$ --- unbounded.
  \item \textbf{Symplectic Euler}: energy error oscillates around a constant
    of order $\mathcal{O}(h)$, bounded for all time.
\end{itemize}
The $\mathcal{O}(h)$ offset between $H$ and $\tilde H$ is why T1 sees an
energy drift of $6\times 10^{-4}$ (bounded, not growing).

% ─────────────────────────────────────────────────────────────────────────────
\section{Exact Discrete Solution: Free Fall}

\subsection{Derivation}

For vertical free fall $H = p^2/(2m) + mgz$, $F = -mg$, the symplectic
Euler update becomes
\begin{align}
  v_{n+1} &= v_n - g h, \label{eq:ff_v}\\
  z_{n+1} &= z_n + h\, v_{n+1} = z_n + h(v_n - gh). \label{eq:ff_z}
\end{align}
where $v_n = p_n/m$.  Telescoping \eqref{eq:ff_v}:
\begin{equation}\label{eq:vn}
  v_n = v_0 - n g h.
\end{equation}
Telescoping \eqref{eq:ff_z}:
\begin{align}
  z_n &= z_0 + h \sum_{k=1}^{n} v_k
       = z_0 + h \sum_{k=1}^{n}(v_0 - k g h) \notag\\
      &= z_0 + n v_0 h - g h^2 \sum_{k=1}^{n} k
       = z_0 + n v_0 h - g h^2 \cdot \frac{n(n+1)}{2}. \label{eq:zn}
\end{align}
Setting $t = nh$ so that $n = t/h$:
\begin{equation}\label{eq:zn_t}
  \boxed{z(t) = z_0 + v_0 t - \tfrac{1}{2} g\, t\,(t + h).}
\end{equation}

\subsection{Comparison with Continuous Theory}

The continuous solution is $z_{\mathrm{cont}}(t) = z_0 + v_0 t -
\frac{1}{2}gt^2$.  Expanding \eqref{eq:zn_t}:
\[
  z(t) = z_{\mathrm{cont}}(t) - \frac{gh}{2}\,t.
\]
The discrete solution lags the continuous one by $\frac{1}{2}ght$, which is
$\mathcal{O}(h)$ at each instant---confirming first-order accuracy.  At
$t = 3\,\mathrm{s}$ and $h = 1/240\,\mathrm{s}$ (PhysX default), the lag is
$\frac{1}{2}\times 9.81 \times \frac{1}{240} \times 3 \approx 0.061\,\mathrm{m}$,
which exceeds the test tolerance of $0.01\,\mathrm{m}$ if the wrong
(continuous) formula is used for comparison.  Using \eqref{eq:zn_t}
reduces the residual to $< 10^{-4}\,\mathrm{m}$ (the test observes
$0.0001\,\mathrm{m}$).

\begin{table}[h]
  \centering
  \begin{tabular}{lcc}
    \toprule
    & Position formula & Area-preserving?\\
    \midrule
    Explicit Euler   & $z_0 + v_0 t - \frac{1}{2}g t^2$ & No\\
    Symplectic Euler & $z_0 + v_0 t - \frac{1}{2}g\,t(t+h)$ & Yes\\
    Continuous       & $z_0 + v_0 t - \frac{1}{2}g t^2$  & (exact)\\
    \bottomrule
  \end{tabular}
  \caption{Free-fall position formulas for $n = t/h$ steps.}
\end{table}

% ─────────────────────────────────────────────────────────────────────────────
\section{Summary}

The chain of equivalences is:
\[
  \underbrace{\text{Type-2 generating function}}_{\text{canonical transformation}}
  \;\Rightarrow\;
  \underbrace{\text{symplectic map}}_{\J^{\top}\Omg\J = \Omg}
  \;\Rightarrow\;
  \underbrace{\text{area-preserving}}_{\det \J = 1}
  \;\Rightarrow\;
  \underbrace{\text{bounded energy error}}_{\tilde H = \text{const}}.
\]
The one-line algorithmic change---replace $p_n$ with $p_{n+1}$ in the
position update---makes $\alpha = \beta\gamma$ cancel, flips an expanding
map to an area-preserving one, and gives PhysX the energy stability it needs
for real-time simulation.

% ─────────────────────────────────────────────────────────────────────────────
\begin{thebibliography}{9}

\bibitem{Hairer2006}
E.~Hairer, C.~Lubich, and G.~Wanner,
\textit{Geometric Numerical Integration: Structure-Preserving Algorithms for
Ordinary Differential Equations}, 2nd~ed.,
Springer Series in Computational Mathematics, vol.~31.
Springer, Berlin, 2006.
Chapter~VI covers symplectic integration; generating functions appear in
\S~VI.5.

\bibitem{Leimkuhler2004}
B.~Leimkuhler and S.~Reich,
\textit{Simulating Hamiltonian Dynamics}.
Cambridge Monographs on Applied and Computational Mathematics, vol.~14.
Cambridge University Press, Cambridge, 2004.

\bibitem{Goldstein2002}
H.~Goldstein, C.~Poole, and J.~Safko,
\textit{Classical Mechanics}, 3rd~ed.
Addison-Wesley, San Francisco, 2002.
Generating functions of canonical transformations: Chapter~9.

\bibitem{Lubich2006chap}
E.~Hairer, G.~Wanner, and C.~Lubich,
``Chapter~VI\@: Symplectic Integration of Hamiltonian Systems,''
lecture notes, Universit\"{a}t T\"{u}bingen, 2006.
Available: \url{https://na.uni-tuebingen.de/~lubich/chap6.pdf}.

\bibitem{Brorson2024}
S.~Brorson,
``Symplectic integrators,''
in \textit{Numerically Solving Ordinary Differential Equations},
Mathematics LibreTexts, 2024.
Available: \url{https://math.libretexts.org/Bookshelves/Differential_Equations/Numerically_Solving_Ordinary_Differential_Equations_(Brorson)/01:_Chapters/1.07:_Symplectic_integrators}.

\end{thebibliography}

\end{document}
